% =============================================================================
%  正文 4：§8 其它群与域、注释、参考文献
% =============================================================================

\section{其它群与域}
\label{sec:other}

\subsection{}

粗略地说，对任何定义在 $\bbQ$ 上的代数群 $G$，Langlands 猜想应该存在一个对应
\begin{equation*}
\left\{\begin{array}{c}
\text{漂亮同态}\\
\Gal(\Qbar/\bbQ)\longrightarrow{}^{L}G(\bbC)
\end{array}\right\}
\;\longrightarrow\;
\left\{\begin{array}{c}
G(\bbA_{\bbQ})\ \text{的漂亮自守表示}\\
\end{array}\right\}
\tag{8.1.1}
\end{equation*}
使得
\begin{equation*}
\{\text{Frobenius 的特征值}\}\;\longrightarrow\;\{\text{Hecke 算子的特征值}\}.
\tag{8.1.2}
\end{equation*}

\subsection{对偶群}
\label{sec:dual}

\subsection{Hermitian 情形}
\label{sec:hermitian}

\subsection{局部域}
\label{sec:localfield}

\subsection{几何 Langlands}
\label{sec:geom}

\subsection{Langlands 函子性猜想}
\label{sec:functoriality}

\subsection{重大结果抽样}
\label{sec:sample}

\begin{yizhu}[原文此处空缺（二）：六节只有标题]
原文第 8.2–8.7 节只有标题，正文全部为空，作者在开头已声明这是 work in progress。
按标题与上下文，这六节本该依次讲：
（8.2）对偶群 ${}^{L}G$ 的构造——Langlands 对偶是 $(G,T)$ 连同根数据作出的复约化群，$G=\GL(n)$ 时 ${}^{L}G=\GL(n,\bbC)$；
（8.3）Hermitian 情形——西群 $U(n)$ 与 unitary 群的朗兰兹对应，对应的是 Hermitian 形式而非一般线性群；
（8.4）局部域——把 $\Gal(\overline{\bbQ}_p/\bbQ_p)$ 与 $G(\bbQ_p)$ 的容许表示对应起来，这是整体猜想的局部影子；
（8.5）几何 Langlands——把 $\Gal$ 换成曲线的基本群、把自守表示换成 $\Bun_G$ 上的层；
（8.6）函子性——群同态 ${}^{L}H\to{}^{L}G$ 应当诱导自守表示的转移，这是朗兰兹猜想的主干推广；
（8.7）重大结果抽样——文中 [F] 那篇综述里列的进展。
以上是译者按标题补的上下文，原文没有任何对应文字，译文不代拟正文。
其中（8.5）几何 Langlands 是本库的落脚点，读法建议见译者说明。
\end{yizhu}

\begin{keytakeaways}
把这一节与第 7 节连起来读，本文的逻辑就完整了：
第 7 节说「$n\ge2$ 的上同调路线走不通」，第 8 节说「所以出路是换群、换域、换范畴」。
8.1 的两行箭头是全篇的最终形式——左边是 Galois 侧，右边是自守侧，中间是「漂亮同态」；
8.5 的几何 Langlands 就是把这两行箭头整体搬到曲线的世界。
本库已译的 SY02-40（Sandvik）、SY02-38（Anschütz）、SY99-04（Raskin），
都是这条箭头在 $n$ 更大、域换成函数域之后的具体化。
\end{keytakeaways}

\bigskip
\section*{注释}
\addcontentsline{toc}{section}{注释}
\begingroup
\small
\begin{enumerate}[label=\arabic*., leftmargin=1.9em, itemsep=1.2pt, topsep=2pt]
\item 换言之，它就是 Cauchy 列的等价类集合；两列 $\{x_n\}$、$\{y_n\}$ 视为等价，若 $|x_n-y_n|_p\to0$。
\item 范数 $|\ |_p$ 定义如下。设 $a/b$ 是既约分数。若 $a$、$b$ 都不被 $p$ 整除，则 $|a/b|_p=1$。
若 $a$ 被 $p$ 整除，写 $a=p^ma'$，则 $|a/b|_p=p^{-m}$。
若 $b$ 被 $p$ 整除，写 $b=p^nb'$，则 $|a/b|_p=p^{n}$。
若把 $c=a/b$ 展开成 $p$ 的幂级数，则 $|a/b|_p=p^{-\val c}$。
\item 这样的嵌入有很多，它们关于 $\bbC$ 上通常的拓扑从不连续。
\item Galois 表示 $\rho$ 总假定关于 $\Gal(\Qbar/\bbQ)$ 的自然拓扑连续，这个拓扑来自它的投射有限群结构；见 [T] \S1。
\item 更一般地，对任何素数 $p\ne\ell$，令 $I_p=\ker(\pi_p)$ 为「惯性群」，$H^{I_p}$ 为 $H$ 中被 $I_p$ 固定的子空间。
若 $H$ 在 $p$ 处非分歧，则 $H^{I_p}=H$。但一般地，商 $\Gal(\Qbar_p/\bbQ_p)/I_p\cong\Gal(\overline{\bbF}_p/\bbF_p)$ 作用在 $H^{I_p}$ 上，
于是得到「Frobenius 元素」$\Fr_p\in\GL(H^{I_p})$。则「$p$ 处」的 $\zeta$ 函数定义为 $1/\det(I-t\Fr_p)$，其中 $I-\Fr_p$ 作用在 $H^{I_p}$ 上。
于是 (4.4.1) 作为所有素数 $p\ne\ell$ 的局部因子之积讲得通。若 $p=\ell$，相应的局部因子牵涉更多技术性考虑。
\item 一般地，若 $G$ 是 Lie 群、$K$ 是闭子群，则 $K$ 的任何表示 $\rho:K\to\GL(V)$（$V$ 为复向量空间）决定 $D=G/K$ 上的齐性向量丛 $E=G\times_KV$，
它由等价类 $[g,v]$ 组成，其中对一切 $k\in K$ 有 $[gk,v]\sim[g,\lambda(k)v]$，并带有 $G$ 的作用 $g\cdot[h,v]=[gh,v]$。
$E$ 的每个（光滑）截面由（光滑）函数 $\tilde f:G\to\bbC$ 给出，满足 $\tilde f(gk)=\lambda(k^{-1})\tilde f(g)$。
$E$ 的（光滑）\textbf{自守因子}是光滑映射 $J:G\times D\to\GL(V)$，满足
\begin{enumerate}[label=(\arabic*), leftmargin=2.2em, topsep=1pt, itemsep=0pt]
\item $J(gg',x)=J(g,g'x)J(g',x)$，对一切 $g,g'\in G$ 与 $x\in D$；
\item $J(k,x_0)=\rho(k)$，对一切 $k\in K$，其中 $x_0\in D$ 是由 $K$ 确定的基点。
\end{enumerate}
自守因子 $J$ 由它在基点的值 $J(g,x_0)$ 决定：任何光滑映射 $j:G\to\GL(V)$ 满足 $j(gk)=j(g)\rho(k)$（对一切 $k\in K$、$g\in G$），
都唯一地延拓为自守因子 $J$，取 $J(g,hx_0)=j(gh)j(h)^{-1}$。
若 $E$ 的自守因子存在，它决定 $E$ 的一个光滑平凡化 $\Phi_J:G\times_KV\to D\times V$，$[g,v]\mapsto(gK,J(g,x_0)v)$，
它关于 $D\times V$ 上的 $J$-自守 $G$ 作用 $g\cdot(x,v)=(gx,J(g,x)v)$ 是 $G$-等变的。
（反之，$E$ 的任何光滑平凡化 $\Psi:E\to D\times V$ 决定唯一的自守因子 $J$ 使 $\Psi=\Phi_J$。）
自守因子使人能把光滑映射 $f:D\to V$ 与 $E$ 的光滑截面 $\tilde f:G\to V$ 等同起来，取 $\tilde f(g)=J(g,x_0)^{-1}f(gK)$。
在某些 $\gamma\in G$ 下不变（即 $\tilde f(\gamma g)=\tilde f(g)$ 对一切 $g\in G$）的截面 $\tilde f$，
就对应于满足 $f(\gamma x)=J(\gamma,x)f(x)$（对一切 $x\in D$）的映射 $f:D\to V$。
在 $\sL(2)$ 的情形，$K=\sO(2)$ 的一维表示
\[
\rho\begin{pmatrix}a&b\\-b&a\end{pmatrix}=(a+ib)^k
\]
有自守因子 $J\!\left(\bigl(\begin{smallmatrix}a&b\\c&d\end{smallmatrix}\bigr),z\right)=(cz+d)^k$，其中 $z\in\he=G/K$。更多细节见 [GP]。
\item Hecke 算子也可以描述为双陪集 $\Gamma g\Gamma$ 对 $f$ 的作用如下。取 $\Gamma'\backslash\Gamma$ 的代表 $\gamma_1,\dots,\gamma_t$，使 $\Gamma=\bigsqcup_{i=1}^{t}\Gamma'\gamma_i$。则对任何 $x\in\he$，
\begin{equation*}
T_g(f)(\Gamma x)=\sum_{i=1}^{t}f(\Gamma g^{-1}\gamma_ix).
\tag{8.7.1}
\end{equation*}
\item $\sL(n)$ 的 Hecke 代数与 $\GL(n)$ 的略有不同。这里有几个例子。
令 $M_n(\bbZ)^+$ 为所有元素为整数、行列式为正的 $n\times n$ 矩阵之环。
则（[Sh1] 定理 3.20）$\sL_n(\bbZ)\backslash M_n(\bbZ)^+/\sL_n(\bbZ)$ 的 Hecke 代数是诸 $H_p$（$p$ 素）之和，每个 $H_p$ 同构于 $n$ 个生成元上的多项式代数
\begin{equation*}
\diag(1,p,\dots,p),\ \diag(1,1,p,\dots,p),\ \dots,\ \diag(p,p,\dots,p).
\tag{8.7.2}
\end{equation*}
$\GL(n,\bbZ)\backslash\GL(n,\bbQ)/\GL(n,\bbZ)$ 的 Hecke 代数是诸 $H_p$ 之和，
每个（[A] 定理 3.2.17）同构于 $n+1$ 个生成元上的多项式代数：即 (8.7.2) 中的那些，再加上 $\diag(p,p,\dots,p)^{-1}$。
\item 其它商 $\Gamma_1\backslash\he$（$\Gamma_1$ 是 $\GL(2,\bbQ)$ 的算术子群）同样可表为
$\GL(2,\bbQ)Z(\bbR)\backslash\GL(2,\bbA)/\sO(2)K^1_{\mathrm f}$，其中 $K^1_{\mathrm f}$ 是有限平代群 $\GL(2,\bbA_{\mathrm f})$ 的适当紧开拓子群。
\item 函数 $\varphi$ 应缓慢增长，或具中等增长，意即对任何 $c>0$ 与任何紧子集 $T\subset\GL(n,\bbA)$，存在常数 $C$、$N$ 使得
\[
\varphi\left(\begin{pmatrix}1&0\\0&a\end{pmatrix}g\right)\le C|a|^N
\]
只要 $a\in\bbA^{\times}$ 且 $|a|>c$；参看注 14。
\item 函数 $\tilde f$ 视为 $\GL(2,\bbR)$ 上的函数，是 Laplace 算子的特征函数，特征值为 $-\frac{k}{2}\left(\frac{k}{2}-1\right)$。
\item $\tilde f$ 定义在商空间 (6.2.1) 上，这意味着 $\tilde f$ 可看成 $\GL(2,\bbA)$ 上的函数，满足 $f(z\gamma g)=f(g)$ 对所有 $z\in Z(\bbA)$ 与 $\gamma\in\GL(2,\bbQ)$。
关于 $K_{\infty}=\sO(2)$ 的等变性表为 $\tilde f(gu)=e^{-ik\theta}f(g)$，其中 $u=\bigl(\begin{smallmatrix}\cos\theta&\sin\theta\\-\sin\theta&\cos\theta\end{smallmatrix}\bigr)$。
\item $\displaystyle\int_{\bbQ\backslash\bbA}\tilde f\left(\begin{pmatrix}1&x\\0&1\end{pmatrix}g\right)dx=0$，对一切 $g\in\GL(2,\bbA)$。
\item 函数 $\varphi$ 应缓慢增长，意即存在常数 $C$、$N$ 使 $|\varphi(g)|\le C\|g\|^N$，其中 $\|\ \|$ 是「高度函数」，[Bu] p.300。
（可以取向量 $(g,\det(g)^{-1})\in\bbA^{2n}\oplus\bbA$ 的长度，用 $\GL(n,\bbA)$ 到全体 $n\times n$ 矩阵空间 $M_n(\bbA)$ 的嵌入。）
若 $\varphi$ 缓慢增长，则它平方可积。
\item 函数 $\varphi$ 应 $z(\mathfrak g)$-有限，意即 $\mathfrak g=\mathfrak{gl}(n,\bbR)$ 的万有包络代数中心 $\mathfrak z$ 中存在有限余维的理想，它把函数 $\varphi$（视为 $\GL(n,\bbR)$ 上的函数）打到零。
\item 函数 $\varphi$ 应 $K$-有限，意即它经 $K$ 平移后张成的函数空间有限维。这里 $K=\sO(n)\times K_{\mathrm f}$，$K_{\mathrm f}=\prod_p\GL(n,\bbZ_p)$ 如前。
\item 更一般地，对任何特征 $\psi:\bbA^{\times}\to\bbC^{\times}$，型为 $\psi$ 的尖点自守形式是函数 $\varphi:\GL(n,\bbQ)\backslash\GL(n,\bbA)\to\bbC$，
满足 $\varphi(zg)=\psi(z)\varphi(g)$ 对所有 $z\in Z(\bbA)\cong\bbA^{\times}$，连同增长、光滑、等变与尖点条件。
\item 群 $\GL(n)$ 的一个特殊性是：每个不可约表示在这个分解中至多出现一次。
\item 实际上 $\GL(n,\bbA)$ 并不作用在 $A_0$ 上，因为 $K$-有限性在平移算子 $R_g$ 下不被保持。
然而 $A_0$ 既是 $(\mathfrak g_{\infty},K_{\infty})$ 模（其中 $\mathfrak g_{\infty}=\mathrm{Lie}(\GL(n,\bbR))$、$K_{\infty}=\sO(n)$），又是 $\GL(n,\bbA_{\mathrm f})$ 模。
作为这样的模它分解为不可约成分，特别地，每个自守表示并不是 $\GL(n,\bbA)$ 的表示，而是
\[
(\mathfrak g_{\infty},K_{\infty})\times\GL(n,\bbA_{\mathrm f})
\]
模。
\item 事实证明 $V^{(K)}$ 由 $\mathfrak z(\mathfrak g_{\infty})$ 的特征向量（或特征函数）组成，相对于某个同态 $\mathfrak z(\mathfrak g_{\infty})\to\bbC$，
而且每个这样的向量还满足注 14 的「缓慢增长」假设；事实上它甚至是速降的。见 [Kn] 定理 7.3。
\item 意思是：对充分大的 $p$，成分 $\pi'_p$ 是一维平凡表示。
\item $\GL(n,\bbQ_p)$ 的表示称为\textbf{容许}的，若它在 $K_p=\GL(n,\bbZ_p)$ 下分解为有限维表示的直和，且每个只出现有限多次。
\item 关于 $G_p$ 上自然定义的 Haar 测度。
\item 若 $h\in H_p$ 是光滑紧支复值 $K_p$-双不变函数，$v\in V^{K_p}$，则
\[
\pi_p(h)(v)=\int_{G_p}h(g)\pi_p(g)v\,dg .
\]
若 $h$ 是单个双陪集 $K_paK_p$ 的特征函数，这给出
\[
\pi_p(h)(v)=\int_{K_p}\pi(k)\pi(a)v\,dk,
\]
它容易被看出仍被 $K_p$ 固定。于是 $\pi_p(h)(v)=\lambda(h)v$，其中 $\lambda(h)\in\bbC$。
容易看出 $\lambda$ 实际上是代数同态 $\lambda:H_p\to\bbC$。
\item 一般地，若 $G$ 是定义在 $\bbQ_p$ 上的约化线性代数群，球 Hecke 代数 $H_G$ 是双陪集 $h\in G(\bbZ_p)\backslash G(\bbQ_p)/G(\bbZ_p)$ 的形式有限线性组合的全体；
等价地，是 $G(\bbQ_p)$ 上局部常值、紧支、$G(\bbZ_p)$-双不变的复值函数 $\varphi$ 连同卷积积的全体。它是交换的。
当 $G=T$ 是分裂环面、余特征群 $X_*(T)=\Hom(\Gm,T)$ 时，同态 $\val:\bbQ_p\to\bbZ$ 延拓为同构
\[
T(\bbQ_p)/T(\bbZ_p)\cong X_*(T),
\]
进而诱导同构 $H_T\cong\bbC[X_*(T)]$，即 Hecke 代数与 $X_*(T)$ 的群代数之间的同构。
令 $T^{\wedge}$ 为对偶环面，$X^*(T^{\wedge})=X_*(T)=\Hom(T,\Gm)$。总之有自然等同
\[
\Hom_{\mathrm{alg}}(H_T,\bbC)\cong\Hom_{\mathrm{alg}}(\bbC[X_*(T)],\bbC)\cong X_*(T)\otimes_{\bbZ}\bbC=X^*(T^{\wedge})\otimes_{\bbZ}\bbC=T^{\wedge}(\bbC),
\]
即群同态 $\bbC[X_*(T)]\to\bbC$ 之群与对偶环面的复点之间。

现在设 $G$ 如上在 $\bbQ_p$ 上分裂。若 $T$ 是 $G$ 的极大环面，则 Satake 变换（见综述 [Gr]）给出同构
$H_G\cong H_T^{W}$，即 $G$ 的 Hecke 代数与 $H_T$ 中 Weyl 群不变元之间的同构，从而诱导等同
\[
\Hom_{\mathrm{alg}}(H_G,\bbC)\cong\Hom_{\mathrm{alg}}(H_T,\bbC)/W\cong T^{\wedge}(\bbC)/W,
\]
即代数同态 $H_G\to\bbC$ 之群与对偶环面 $T^{\wedge}(\bbC)$ 复点中 Weyl 群轨道之间。
集合 $T^{\wedge}(\bbC)/W$ 也参数化对偶群 $G^{\wedge}$ 的复点群中的半单共轭类。
当 $G=\GL(n)$ 时，极大环面 $T$ 与对偶环面 $T^{\wedge}$ 可等同，$G$ 与 $G^{\wedge}$ 亦可等同。
\item 这些数本质上是对应「古典」模形式 $f$ 上 $\diag(1,\dots,1,p,\dots,p)$ 作用之特征值的基本对称函数；见 [K]。
\item 有定理保证：除有限多个素数 $p$ 外，$\pi_p$ 都是非分歧的。然而 $\pi$ 的 $L$ 函数应包含所有素数的贡献，对分歧的素数必须另作处理。
\item 用现代语言，这是 $X$ 的 Baily–Borel 紧化 $\overline X$ 的相交上同调
\[
IH^1(\overline X,E_k;\bbQ)\cong H^1_P(X,E_k;\bbQ).
\]
系数在 $E_k$ 中的相交上同调讲得通，即使局部系统 $E_k$ 可能无法延拓到 $X$ 的 Baily–Borel 紧化上。
相交上同调也有平展版本与比较同构
\[
IH^i_{\et}(\overline X,E_k;\bbQ_{\ell})\cong IH^i(\overline X,E_k;\bbQ)\otimes_{\bbQ}\bbQ_{\ell},
\]
介于 $\overline X$（视为定义在有理数上的代数簇）的平展相交上同调与 $\overline X$ 的（拓扑）相交上同调之间。
\end{enumerate}
\endgroup

\clearpage
\section*{参考文献}
\addcontentsline{toc}{section}{参考文献}
\markboth{参考文献}{参考文献}
\begingroup
\small
\begin{list}{}{\leftmargin=2.4em\labelwidth=2.0em\labelsep=0.4em\itemsep=1.5pt}
\item[[A]] A.\,N.\ Andrianov, \emph{Quadratic Forms and Hecke Operators}, Grundlehren der mathematischen Wissenschaften 286, Springer Verlag, Berlin, 1987.
\item[[Ar]] J.\ Arthur, The principle of functoriality, \emph{Bull.\ Amer.\ Math.\ Soc.}\ 40 (2002), 39--53.
\item[[Ba]] T.\ Bailey and A.\ Knapp, ed., \emph{Representation Theory and Automorphic Forms}, Proc.\ Symp.\ Pure Math.\ 61, Amer.\ Math.\ Soc., Providence, R.\,I., 1977.
\item[[Be]] J.\ Bernstein and S.\ Gelbart, ed., \emph{An Introduction to the Langlands Program}, Birkh\"auser, Boston MA, 2003.
\item[[Bu]] D.\ Bump, \emph{Automorphic Forms and Representations}, Cambridge Studies in Advanced Mathematics 55 (1997), Cambridge University Press, New York.
\item[[BR]] D.\ Blasius and J.\ Rogawski, Zeta functions of Shimura varieties, in \emph{Motives}, Proc.\ Symp.\ Pure Math.\ 55 (part 2), 525--572, Amer.\ Math.\ Soc., Providence R.\,I., 1994.
\item[[F]] E.\ Frenkel, Recent advances in the Langlands program, \emph{Bull.\ Amer.\ Math.\ Soc.}\ 41 (2004), 151--184.
\item[[G1]] S.\ Gelbart, \emph{Automorphic Forms on Adele Groups}, Annals of Math Studies 83, Princeton University Press, Princeton NJ, 1975.
\item[[G2]] S.\ Gelbart, An elementary introduction to the Langlands Program, \emph{Bull.\ Amer.\ Math.\ Soc.}\ 10 (1984), 177--219.
\item[[GF]] I.\,M.\ Gelfand and S.\ Fomin, Geodesic flows on manifolds of constant negative curvature, \emph{Uspekhi Mat.\ Nauk}\ 7 (1952), 118--137; English translation in \emph{Amer.\ Math.\ Soc.\ Transl.}\ 1 (1965), 49--65.
\item[[GP]] M.\ Goresky and W.\ Pardon, Chern classes of automorphic vector bundles, \emph{Inv.\ Math.}\ 147 (2002), 561--612.
\item[[Gr]] B.\ Gross, On the Satake isomorphism, in \emph{Galois Representations in Arithmetic Algebraic Geometry}, A.\ Scholl and R.\ Taylor, ed., Cambridge University Press, New York, 1998, 223--237.
\item[[JL]] H.\ Jacquet and R.\ Langlands, \emph{Automorphic Forms on GL(2)}, Lecture Notes in Mathematics 114, Springer Verlag, New York, 1970.
\item[[Kn]] A.\ Knapp, Introduction to the Langlands Program, \emph{Proc.\ Symp.\ Pure Math.}\ 61 (1997), 245--302.
\item[[K]] S.\ Kudla, From modular forms to automorphic representations, in \emph{An Introduction to the Langlands Program}, J.\ Bernstein and S.\ Gelbart, ed., Birkh\"auser, Boston MA, 2003, 133--151.
\item[[M]] J.\ Milne, review of book by M.\ Harris and R.\ Taylor, \emph{Math Reviews} 2002m:11050.
\item[[R1]] J.\ Rogawski, The nonabelian reciprocity law for local fields, \emph{Notices of the A.\,M.\,S.}\ 47 (2000), 35--41.
\item[[R2]] J.\ Rogawski, review of books by Borel, Bump, and Iwaniec, \emph{Bull.\ A.\,M.\,S.}\ 35 (1998), 253--263.
\item[[R3]] J.\ Rogawski, review of books by Laumon, \emph{Math Reviews} MR 98C:11045a, 98C:11045b.
\item[[Se]] J.-P.\ Serre, \emph{A Course in Arithmetic}, Graduate Texts in Mathematics 7, Springer Verlag, New York, 1973.
\item[[Sh1]] G.\ Shimura, \emph{Introduction to the Arithmetic Theory of Automorphic Forms}, Math.\ Soc.\ Jap., 1971, reprinted Princeton University Press, Princeton NJ, 1994.
\item[[T]] R.\ Taylor, Galois representations, \emph{Proceedings of the International Congress of Mathematicians, Beijing, 2002}, vol.\,I, World Scientific, New York, 449--474. See also longer version, to appear in \emph{Annales de la Facult\'e des Sciences de Toulouse}.
\end{list}
\endgroup
